\section{Mean-payoff input and arena restrictions}
\label{sec:arenas}

We first record what is needed from ordinary mean-payoff games and
how winning regions behave under attractor removal. These are the
interfaces between the quantitative objective and the recursion.

A finite path $\rho=v_0\cdots v_t$ has length $t$ and weight
\[
                  \wt(\rho)=\sum_{i=0}^{t-1}r(v_i).
\]
Thus assigning edge weight $w(v,u)=r(v)$ preserves every finite
prefix weight and the infinite-play payoff. For an arena $H$, put
$W_H=\max_{v\in H}|r(v)|$ when $H$ is nonempty.
The mean-payoff objective and both winning objectives in
\eqref{eq:objective} are prefix independent: a finite initial path
does not change the recurring priorities or the lower limiting
average. A strategy may therefore be restarted at the current vertex
with its memory initialized there. All statements about such a
restart continue to quantify over arbitrary history-dependent
opposing strategies.

\subsection{The ordinary mean-payoff interface}

\begin{proposition}[Mean-payoff input]\label{prop:mp-interface}
For a finite nonempty total arena $H$ with signed binary integer
weights, its ordinary mean-payoff winning set
\[
 U_{\mathrm{mp}}(H)=
 \{v:\text{Max can force liminf mean payoff at least }0\}
\]
is computable deterministically in
$2^{O((\log_2(L_H+2))^2)}$ bit operations, where $L_H$ is its full
binary encoding length. There is a single positional Max strategy
winning from every vertex of $U_{\mathrm{mp}}(H)$, and a single
positional Min strategy forcing strictly negative liminf mean payoff
from every vertex outside it. Both guarantees hold against every
history-dependent opposing strategy.
\end{proposition}

\begin{proof}
The winning-set bound is
\cite[Theorem~1.1]{OpenAI2026Deterministic}.
Its Corollary~7.1 supplies exact values and globally optimal
positional strategies. Positional attainment identifies
$U_{\mathrm{mp}}(H)$ with the vertices of nonnegative value; at every
other vertex Min's strategy bounds the payoff by its strictly
negative value. These assertions use the same liminf convention as
\eqref{eq:objective}.
\end{proof}

\begin{lemma}[Uniform prefix bound]\label{lem:mp-prefix}
If Max wins ordinary mean payoff from every vertex of a nonempty
arena $H$ of size $N$, then a single positional strategy guarantees
\[
                  \wt(\rho)\ge -(N-1)W_H
\]
for every finite consistent path $\rho$, from every start.
\end{lemma}

\begin{proof}
Fix the global positional strategy of
Proposition~\ref{prop:mp-interface}, retaining its chosen edge at
each Max vertex and all edges at Min vertices. Every simple directed
cycle in this graph has nonnegative weight: otherwise Min could
repeat a negative cycle starting on it, contradicting the global
guarantee. Delete simple cycles successively from any finite walk.
Their nonnegative weights leave a path of at most $N-1$ edges whose
weight is no greater than that of the original walk. Its weight is
at least $-(N-1)W_H$.
\end{proof}

The same cycle argument underlies the winning-policy construction in
\cite[Proposition~2.2]{OpenAI2026Deterministic}.
Only the winning set in Proposition~\ref{prop:mp-interface} is
queried by our algorithm; the strategies are used in its correctness
proof.

\subsection{Attractors and dominions}

We use the standard attractor and dominion properties underlying
reduced-precision parity recursion; compare
\cite[Lemmas~2.1--2.2]{LehtinenParysScheweWojtczak2022}.
Their short proofs below apply to the present prefix-independent
objectives as well.

An \emph{arena restriction} is the induced game on a vertex set that
retains at least one outgoing edge at every retained vertex. We allow
the empty restriction and identify an arena with its vertex set.
A restriction $R$ of $H$ is \emph{closed at player $P$'s vertices}
if every successor in $H$ of a $P$-vertex in $R$ also lies in $R$.
Write $Q$ for the other player.

For $X\subseteq H$, the attractor $\Attr_P^H(X)$ is the least set
containing $X$ that is closed under the following additions: a
$P$-vertex with some successor already in the set, and a $Q$-vertex
with all successors already in the set. Assigning each newly added
vertex its stage of entry gives a strategy that decreases this rank
until reaching $X$. Thus $P$ can force a visit to $X$ from its
attractor within at most $|H|$ edges.

The complement $H\setminus\Attr_P^H(X)$ is a total restriction
closed at $P$'s vertices. Moreover, a total restriction disjoint from
$X$ and closed at $P$'s vertices is disjoint from the attractor:
inductively, no retained $P$-vertex has a successor in the growing
attractor, and each retained $Q$-vertex has a retained successor
outside it.

A \emph{$P$-dominion in $H$} is a set $D$ from each of whose vertices
$P$ can win the given objective while keeping play in $D$, against
every opposing strategy. Such a set is a total restriction closed
at $Q$'s vertices. Conversely, a total restriction with this closure
is a $P$-dominion exactly when $P$ wins everywhere in its induced
game. The empty set is a dominion by convention.

The whole winning region of either player is a dominion for that
player, even before determinacy is known. Along a winning strategy,
the continuation after any consistent finite history still wins,
by prefix independence. Every vertex reached therefore belongs to
the same player's winning region. This also shows that the region
is closed at the opposing player's vertices.

\begin{lemma}[Restriction facts]\label{lem:arena-facts}
Let $P,Q$ be opposite players in an arena $H$.
\begin{enumerate}
\renewcommand{\labelenumi}{\textup{(\roman{enumi})}}
 \item If $D$ is a $P$-dominion in $H$ and $R$ is a restriction
 closed at $P$'s vertices, then $D\cap R$ is a $P$-dominion in $R$.
 \item If $R$ is closed at $P$'s vertices, every $Q$-dominion in
 $R$ is a $Q$-dominion in $H$.
 \item If $D$ is a $P$-dominion in $H$ and $D\cap X=\varnothing$,
 then removing $\Attr_Q^H(X)$ leaves $D$ intact and still a
 $P$-dominion in the residual arena.
\end{enumerate}
\end{lemma}

\begin{proof}
For \textup{(i)}, use a strategy keeping play in $D$. At $P$-vertices
its choices remain in $R$ by closure. At $Q$-vertices every ambient
successor stays in $D$, and the restricted opponent must choose one
in $R$. Hence play remains in the intersection and wins. This also
establishes that the intersection is total.

For \textup{(ii)}, $Q$ can choose the edges of its restricted
dominion strategy, and $P$ cannot exit $R$. The same plays and
winning condition result in $H$.

For \textup{(iii)}, $D$ is closed at $Q$'s vertices and is disjoint
from $X$. The attractor-avoidance observation above shows that it is
disjoint from $\Attr_Q^H(X)$. Its winning strategies and the opposing
choices within $D$ are therefore retained.
\end{proof}
